\SourceLeaf{98}%
imaginary, by an infinite plane, in which each point, determined by abscissa $=a$ and ordinate $=b$, represents, as it were, the quantity $a+bi$. The continuous passage from one value of $x$ to another, $a+bi$, thus takes place along a line and is therefore possible in infinitely many ways. I now assert that the integral $\int \varphi x\mathord{.} dx$ always takes the same value along two different paths if nowhere in the region enclosed by the lines representing the two paths does one have $\varphi x = \infty$. This is a very beautiful theorem\HistoricalNote{\textsuperscript{*})}{*)}{More precisely, it is further assumed here that $\varphi x$ is itself a single-valued function of $x$, or at least that only one system of its values is taken throughout that entire region without any break in continuity.}, whose proof is not particularly difficult and which I shall give on a suitable occasion. It is connected with other beautiful truths concerning expansions in series. The passage to any point can always be made without touching any such place where $\varphi x=\infty$. I therefore require that one avoid such points, where the original conception of $\int \varphi x\mathord{.} dx$ plainly loses its clarity and easily leads to contradictions. This also makes clear how a function generated by $\int \varphi x\mathord{.} dx$ can have several values for the same value of $x$. In passing to that value one can go around a point where $\varphi x=\infty$ either not at all, once, or several times. If, for\,example, one defines $\log x$ by $\int \frac{1}{x}\,dx$, starting from $x=1$, one obtains $\log x$ either without enclosing the point $x=0$ or by going around it once or several times. Each time the constant $+2\pi i$ or $-2\pi i$ is added. This makes the multiple logarithms of every number quite clear. If $\varphi x$ can never become infinite for any finite value of $x$, the integral is always a single-valued function. This is the case, for\,example, for $\varphi x = \dfrac{e^x-1}{x}$, so that $\int \dfrac{e^x-1}{x}\,dx$ is certainly a single-valued function of $x$. Its value is represented by the \textls{always} convergent series, which thus always has one and only \textls{one} meaning,
\[
x+\tfrac{1}{4}xx+\tfrac{1}{18}x^3+\tfrac{1}{96}x^4+\text{etc.}
\]
\noindent I wish that, since Mr \textsc{Soldner} meant to%
